\documentclass[10pt]{article}
\usepackage[T1]{fontenc}
\usepackage[utf8]{inputenc}
\usepackage[french]{babel}
\usepackage{amsmath,amssymb,mathrsfs,mathpazo}
\usepackage{geometry}
\geometry{a4paper,margin=2cm}
\usepackage{graphicx,array,tabularx,tabularray,colortbl,multirow,multicol,booktabs,enumitem,pifont,ulem,fancybox,diagbox,cancel}
\usepackage{pstricks,pstricks-add,pst-plot,pst-tree,pst-node,pst-text,pst-func,pst-math,pst-bspline,pst-3dplot,pst-eucl,pst-solides3d}
\let\pscalculate\undefined
\usepackage{pst-bezier}
\usepackage{tikz,pgfplots,tkz-tab}
\usetikzlibrary{arrows,arrows.meta,patterns,calc,positioning,trees,decorations.markings,shapes.geometric,shapes.misc,shapes.arrows,intersections,decorations.pathmorphing,decorations.pathreplacing,backgrounds,shadows,babel}
\pgfplotsset{compat=1.18}
\usepgfplotslibrary{fillbetween}
\usepackage[np]{numprint}
\usepackage{listings,makecell,pas-tableur,twoopt}
\let\barre\undefined
\usepackage[pstricks]{tablvar}
\definecolor{darkgreen}{RGB}{0,100,0}
\providecommand{\red}{\color{red}}
\providecommand{\blue}{\color{blue}}
\providecommand{\green}{\color{green}}
\providecommand{\black}{\color{black}}
\providecommand{\white}{\color{white}}
\providecommand{\gray}{\color{gray}}
\providecommand{\cyan}{\color{cyan}}
\providecommand{\magenta}{\color{magenta}}
\providecommand{\yellow}{\color{yellow}}
\providecommand{\violet}{\color{violet}}
\providecommand{\orange}{\color{orange}}
\providecommand{\euro}{\texteuro}
\providecommand{\EUR}{\texteuro}
\providecommand{\ieme}{\textsuperscript{e}}
\providecommand{\degres}{\ensuremath{^\circ}}
\providecommand{\degree}{\ensuremath{^\circ}}
\providecommand{\vv}[1]{\overrightarrow{#1}}
\providecommand{\Coord}[1]{#1}
\providecommand{\decosix}{}
\providecommand{\starredbullet}{\ensuremath{\bullet}}
\pagestyle{empty}
\setlength{\parindent}{0pt}
\usepackage{adjustbox}\DeclareUnicodeCharacter{2236}{:}\DeclareUnicodeCharacter{2190}{\ensuremath{\leftarrow}}
\providecommand{\R}{\mathbb{R}}
\providecommand{\N}{\mathbb{N}}
\providecommand{\Z}{\mathbb{Z}}
\providecommand{\Q}{\mathbb{Q}}
\providecommand{\D}{\mathbb{D}}
\providecommand{\C}{\mathbb{C}}
\providecommand{\e}{{\mathrm{e}}}
\providecommand{\vect}[1]{\overrightarrow{#1}}
\providecommand{\vectt}[1]{\overrightarrow{\mathrm{#1}}}
\providecommand{\vvt}[1]{\overrightarrow{\mathrm{#1}}}
\providecommand{\np}[2][]{{#2~#1}}
\providecommand{\vv}[1]{{\overrightarrow{#1}}}
\renewcommand{\d}{{\mathrm{d}}}
\providecommand{\degres}{{^\circ}}
\providecommand{\degre}{{^\circ}}
\providecommand{\blue}{\color{blue}}
\providecommand{\red}{\color{red}}
\providecommand{\enskip}{\quad}
\providecommand{\llbracket}{\lbrack\!\lbrack}
\providecommand{\rrbracket}{\rbrack\!\rbrack}
\providecommand{\euro}{{\text{€}}}
\providecommand{\no}{{\text{n°}}}
\providecommand{\interval}[2]{{[#1;#2]}}
\providecommand{\intervFO}[2]{{[#1;#2[}}
\providecommand{\dotfill}{\ldots}

\providecommand{\ds}{\displaystyle}
\providecommand{\textding}[1]{#1}
\providecommand{\Oij}{\(\mathrm{O};\vec{\imath},\vec{\jmath}\)}
\providecommand{\Oijk}{\(\mathrm{O};\vec{\imath},\vec{\jmath},\vec{k}\)}
\providecommand{\Ouv}{\(\mathrm{O};\vec{u},\vec{v}\)}
\providecommand{\intervFF}[2]{\ensuremath{[#1;#2]}}
\renewcommand{\euro}{\ifmmode\text{\texteuro}\else\texteuro\fi}
\DeclareUnicodeCharacter{00B0}{\ifmmode^{\circ}\else\textdegree\fi}
\DeclareUnicodeCharacter{20AC}{\ifmmode\text{\texteuro}\else\texteuro\fi}
% Énoncés : Éducation nationale. Transcription et corrigés : APMEP, Denis Vergès et contributeurs. Découpage et classement : CoopMaths, Sébastien Lozano, Éric Elter et contributeurs. Adaptation de la présentation : Maths Explorer.
% CC BY-SA : https://coopmaths.fr/www/about/
% https://coopmaths.fr/alea/static/bac/2023/tex/bac_2023_03_sujet1_polynesie_4.tex
\begin{document}
\noindent\textit{APMEP / CoopMaths. CC BY-SA. Présentation adaptée par Maths Explorer.}\par\medskip
\medskip

Soit $\left(u_n\right)$ la suite définie par $u_0 = - 1$ et, pour tout entier naturel $n$ :
\[u_{n+1} = 0,9u_n - 0,3.\]

\begin{enumerate}
\item 
	\begin{enumerate}
		\item Démontrer par récurrence que, pour tout $n \in \N$,\: $u_n  = 2 \times 0,9^n - 3$.
		\item En déduire que pour tout $n \in \N,\: \: - 3 < u_n \leqslant - 1$.
		\item Démontrer que la suite $\left(u_n\right)$  est strictement décroissante.
		\item Démontrer que la suite $\left(u_n\right)$ converge et préciser sa limite.
	\end{enumerate}
\item On se propose d'étudier la fonction $g$ définie sur $]-3~;~-1]$ par : 
\[g(x) = \ln (0,5 x + 1,5) - x.\]

	\begin{enumerate}
		\item Justifier toutes les informations données par le tableau de variations de la fonction $g$ (limites, variations, image de $-1$)
		
\begin{center}
\psset{unit=1cm,arrowsize=2pt 4}
\begin{pspicture}(8.5,3.5)
\psframe(8.5,3.5)\psline(0,3)(8.5,3)\psline(2.5,0)(2.5,3.5)
\uput[u](1.25,2.9){$x$} \uput[u](2.7,2.9){$-3$} \uput[u](5.5,2.9){$-2$}\uput[u](8.2,2.9){$-1$} \uput[d](5.5,3){$g(-2)$}\uput[u](3.15,0){$- \infty$}\uput[u](8.3,0){$1$}
\rput(1.25,1.5){Variations de $g$}
\psline(2.65,0)(2.65,3)\psline(2.75,0)(2.75,3)
\psline{->}(3.2,0.5)(5.2,2.5)\psline{->}(5.8,2.5)(8.2,0.5)
\end{pspicture}
\end{center}
		\item En déduire que l'équation $g(x) = 0$ a exactement une solution que l'on
notera $\alpha$ et dont on donnera un encadrement d'amplitude $10^{-3}$.
	\end{enumerate}
\item Dans la suite de l'exercice, on considère la suite $\left(v_n\right)$ définie pour tout $ \in \N$, par :
\[v_n = \ln \left(0,5 u_n + 1,5\right).\]
	\begin{enumerate}
		\item En utilisant la formule donnée à la question \textbf{1. a.}, démontrer que la suite v est arithmétique de raison $\ln (0,9)$.
		\item Soit $n$ un entier naturel.
		
Démontrer que $u_n = v_n$ si, et seulement si $g\left(u_n\right) = 0$.
		\item Démontrer qu'il n'existe aucun rang $k \in \N$ pour lequel $u_k = \alpha$.
		\item En déduire qu'il n'existe aucun rang $k \in \N$ pour lequel $v_k = u_k$.
	\end{enumerate}
\end{enumerate}

\end{document}
